% ============================================================
% CHAPITRE 10
% ============================================================
\chapter{Tests et qualité logicielle}
\label{ch:qualite}

\section{Trois niveaux de vérification}
La plateforme est vérifiée à trois niveaux (figure~\ref{fig:pyramide}). Les
tests unitaires en Python pur s'exécutent en une seconde et couvrent la
logique métier ; les tests Spark démarrent une vraie session avec GraphFrames
et Delta Lake pour vérifier les motifs et le schéma ; les mesures de bout en
bout (évaluation, ablation, benchmark, test de panne) vérifient la plateforme
entière sur des données réalistes.

\begin{figure}[htp]
\centering
\begin{tikzpicture}
  \fill[vertcode!80] (-4.6,0) -- (4.6,0) -- (3.45,1.15) -- (-3.45,1.15) -- cycle;
  \fill[bleuprojet!80] (-3.4,1.25) -- (3.4,1.25) -- (2.25,2.4) -- (-2.25,2.4) -- cycle;
  \fill[orangeprojet!85] (-2.2,2.5) -- (2.2,2.5) -- (1.1,3.65) -- (-1.1,3.65) -- cycle;
  \node[font=\sffamily\scriptsize\bfseries, text=white, align=center] at (0,0.57) {19 tests unitaires (Python pur)\\[-1pt]{\tiny\mdseries scénarios, score, index glissant, vélocité, identifiants}};
  \node[font=\sffamily\scriptsize\bfseries, text=white, align=center] at (0,1.82) {10 tests Spark\\[-1pt]{\tiny\mdseries motifs, SCC, CC, mules, schéma Delta}};
  \node[font=\sffamily\tiny\bfseries, text=white, align=center] at (0,3.02) {bout en bout\\[-1pt]{\tiny\mdseries 4 mesures}};
  \foreach \y/\t in {0.57/{$\approx$\,1~s · sur l'hôte comme en CI}, 1.82/{$\approx$\,2~min · dans l'image Docker (JVM)}, 3.07/{30--45~min · sur la plateforme complète}}
    \node[font=\sffamily\scriptsize, anchor=west, color=encre2] at (4.9,\y) {\t};
\end{tikzpicture}
\caption{Les trois niveaux de vérification : 29 tests automatisés et quatre mesures de bout en bout.}
\label{fig:pyramide}
\end{figure}

\begin{table}[htp]
\centering\small
\begin{tabular}{@{}>{\raggedright\arraybackslash}p{6.9cm}>{\raggedright\arraybackslash}p{7.3cm}@{}}
\toprule
\sffamily\bfseries Test (extraits) & \sffamily\bfseries Ce qu'il garantit\\
\midrule
\fichier{test\_triangle\_temporal\_unique\_rotation} & un carrousel n'est signalé qu'une fois, pas trois\\
\fichier{test\_triangle\_out\_of\_order\_rejected\ldots} & une boucle non ordonnée dans le temps est écartée\\
\fichier{test\_incremental\_search\_only\_reports\ldots} & le chemin rapide ne redétecte pas les anciens cycles\\
\fichier{test\_long\_cycle\_found\_inside\_scc} & un cycle de 5 comptes est trouvé dans sa SCC\\
\fichier{test\_mule\_velocity\_spike\_excludes\ldots} & un commerçant régulier n'est pas une mule\\
\fichier{test\_alert\_id\_is\_deterministic\ldots} & même fraude, même identifiant (exactly-once)\\
\fichier{test\_delta\_schema\_and\_constraints} & une ligne invalide est refusée par Delta\\
\fichier{test\_every\_scenario\_builds} (5 typologies) & chaque fraude est générable et cohérente\\
\bottomrule
\end{tabular}
\caption{Extraits des tests automatisés (\fichier{tests/}).}
\label{tab:tests}
\end{table}

\section{Intégration continue}
À chaque modification, un pipeline GitHub Actions rejoue la chaîne de
vérification (figure~\ref{fig:ci}) : le lint (\fichier{ruff}), les tests en
Python pur, les tests Spark avec Java~17, puis la génération du DDL Delta du
livrable~3.

\begin{figure}[htp]
\centering
\begin{tikzpicture}[
  s/.style={draw=#1, line width=0.9pt, fill=#1!10, rounded corners=3pt, minimum width=2.55cm,
    minimum height=1.0cm, align=center, font=\sffamily\scriptsize}]
  \foreach \t/\d/\c [count=\i] in {
    {Dépendances}/{\fichier{requirements.txt}}/grismoyen,
    {Lint}/{\fichier{ruff check .}}/bleuprojet,
    {Tests Python}/{\fichier{pytest -m "not spark"}}/bleuprojet,
    {Tests Spark}/{Java 17, GraphFrames, Delta}/orangeprojet,
    {DDL Delta}/{\fichier{--print-ddl}}/vertcode} {
    \node[s=\c] (c\i) at (\i*3.05,0) {\textbf{\t}\\{\tiny \d}};
  }
  \foreach \a/\b in {1/2,2/3,3/4,4/5} \draw[fleche] (c\a) -- (c\b);
\end{tikzpicture}
\caption{Le pipeline d'intégration continue (\fichier{.github/workflows/ci.yml}).}
\label{fig:ci}
\end{figure}

Au-delà des tests, le job écrit à chaque micro-batch et à chaque passage du
chemin lent une ligne de métriques (durées par étape, lignes traitées, taille
du graphe, retards écartés). Ces journaux ont été le premier outil de
diagnostic de chaque problème de performance du chapitre~\ref{ch:performance}
et ont servi à tracer la plupart des figures de ce rapport.
